% ============================================================================
% Fichier  : annexes/Ann_E_bibliographie.tex
% Version  : 1.0 -- novembre 2026
% ============================================================================

\chapter{Bibliographie Commentée}
\label{ann:bibliographie}

\begin{description}
    \item[Les textes.] La loi n°~2024/017~\cite{loi2024017} est la source
          première du cours~; on la lit en entier, une fois, avant
          d'auditer. La loi n°~2010/012~\cite{loi2010012} reste utile pour
          la vie privée dans les communications.
    \item[Le cadre africain.] La Convention de Malabo~\cite{ua2014malabo}
          et l'acte additionnel de la CEDEAO~\cite{cedeao2010acte} situent
          la loi camerounaise dans son environnement régional.
    \item[La comparaison européenne.] Le RGPD~\cite{ue2016rgpd} et les
          recommandations du Comité européen~\cite{edpb2020transferts}
          servent de référence, à condition de ne pas les confondre avec la
          loi camerounaise.
    \item[Les normes.] L'ISO/IEC~29100~\cite{iso29100} pour le vocabulaire,
          l'ISO/IEC~27701~\cite{iso27701} pour le système de management,
          l'ISO/IEC~29134~\cite{iso29134} pour l'évaluation d'impact,
          l'ISO~19011~\cite{iso19011_2018} pour la conduite de l'audit.
    \item[La méthode PIA de la CNIL.]~\cite{cnil2018pia} La plus accessible
          des méthodes d'évaluation d'impact~; ses échelles sont reprises au
          Chapitre~\ref{chap:risque}.
    \item[La vie privée.] D.~J.~Solove~\cite{solove2008} et
          H.~Nissenbaum~\cite{nissenbaum2010}, pour comprendre ce que la
          protection des données protège.
    \item[La recherche de l'auteur.] « Vers l'auditologie »~\cite{minka2026auditologie}
          et l'instrument trichotomique~\cite{minka2026nek}~; le manuscrit
          en préparation sur la chaîne pénale~\cite{minka2026chainepenale}.
\end{description}
